This study evaluates an evidence-constrained assistance workflow, not an autonomous fault classifier. Its principal limits concern the interpretation of the signal observations, the coverage of retrieval evaluation, and the scope of the embedding comparison.

\subsection{Absence of ground truth and observation-to-fault inference}
The raw measurement CSV records used in this study do not provide a uniformly aligned, gold-standard fault taxonomy. The deterministic probes therefore report signal events---for example, spikes, zero values, extreme values, and cross-channel co-movement---rather than confirmed equipment faults. Generated condition reports are derived workflow artifacts, not ground-truth annotations. They demonstrate traceability from data to an observation record, but cannot establish end-to-end diagnostic accuracy or root-cause accuracy. Engineering review and, where available, maintenance records remain necessary to determine the physical mechanism behind an observed event.

\subsection{Retrieval corpus and relevance coverage}
The robustness study is bounded by a fixed 64-document corpus, consisting of a 13-document IAEA core and 51 same-domain supplementary documents. Its corpus-expansion protocol is useful for testing whether designated authoritative sources remain visible, but it is not a representative sample of every plant documentation environment. The scope experiments also use a small, fixed bilingual query set and do not include a human relevance judgment for each retrieved passage or an expert audit of every citation-bearing memo claim. Consequently, strict chunk and document-level self-retrieval are retrieval diagnostics, not direct measures of clinical or operational usefulness.

\subsection{Embedding and generation boundaries}
The historical index results and the new CPU comparison use different chunk constructions and are therefore reported separately. The latter compares two MiniLM encoders under one fixed retrieval configuration; it establishes that encoder choice affects bilingual similarity scales and source-priority decay, but it does not rank all available embedding models. The study also evaluates retrieval rather than generation-side factuality, completeness, or calibration. Any generated diagnostic memo should remain linked to retrieved passages and be reviewed by a qualified engineer before it informs a safety-relevant decision.

\subsection{Threats to validity and next validation stage}
The reported relationships between corpus expansion, cosine similarity, and source priority may change with the corpus, metadata policy, chunking implementation, query formulation, encoder, or reranker. Repeated measurements reduce the sensitivity of the reported curves to random supplementary-document selections, but they do not eliminate these external-validity limits. A next validation stage should separately assess observation correctness against reviewed cases, passage-level support for each memo claim, expert usefulness of the memo, and generation-side factuality under a preregistered evaluation protocol.
